\chapter{METODOLOGIA}

\orientacao{Descreva COMO você pretende conduzir a pesquisa no TCC II. Inclua: tipo de pesquisa, fontes de dados, procedimentos, ferramentas e critérios de análise. Seja específico o suficiente para que outro pesquisador possa replicar seu estudo.}

\placeholder{Parágrafo introdutório: classifique a pesquisa quanto à natureza (ex.: aplicada, básica), abordagem (qualitativa, quantitativa ou mista), objetivos (exploratória, descritiva, explicativa) e procedimento (pesquisa bibliográfica, estudo de caso, experimento, survey, etc.). Justifique cada escolha.}

\section{Tipo e classificação da pesquisa}

\orientacao{Use a classificação de Gil (2010) ou similar. Exemplo: ``Esta pesquisa é de natureza aplicada, abordagem quantitativa, objetivos descritivos e procedimento bibliográfico-experimental.''}

\placeholder{Descreva e justifique a classificação adotada.}

\section{Fonte e coleta de dados}

\orientacao{Indique de onde virão os dados: bases públicas, APIs, questionários, entrevistas, repositórios, simulações, etc. Se usar dados secundários, cite a fonte oficial.}

\placeholder{Descreva a(s) fonte(s) de dados. Exemplo: ``Os dados serão obtidos de [fonte], disponível em [URL ou referência], referente ao período [X].''}

\placeholder{Se aplicável, liste os arquivos, variáveis ou instrumentos de coleta que serão utilizados.}

\section{Amostra ou escopo}

\orientacao{Defina critérios de inclusão/exclusão. Para pesquisas quantitativas, descreva a população e a amostra. Para qualitativas, descreva os participantes ou casos.}

\placeholder{Descreva o escopo da amostra ou delimitação do estudo. Exemplo: ``Serão considerados [critério 1], excluindo [critério 2].''}

\placeholder{Se já souber o tamanho estimado da amostra, informe aqui. Caso contrário, indique como será determinado.}

\section{Variáveis e instrumentos}

\orientacao{Liste variáveis independentes, dependentes e de controle (se houver). Descreva instrumentos: questionários, scripts, ferramentas de análise.}

\begin{itemize}
    \item \textbf{Variável independente:} \placeholder{nome e descrição};
    \item \textbf{Variável dependente:} \placeholder{nome e descrição};
    \item \textbf{Variáveis de controle (opcional):} \placeholder{nome e descrição}.
\end{itemize}

\section{Procedimentos de análise}

\orientacao{Detalhe as etapas analíticas: tratamento de dados, testes estatísticos, técnicas de ML, análise de conteúdo, etc. Se usar testes estatísticos, formalize hipóteses (H$_0$ e H$_1$).}

\placeholder{Descreva os procedimentos passo a passo. Exemplo: (1) limpeza dos dados; (2) análise descritiva; (3) testes inferenciais; (4) visualização.}

\placeholder{Informe ferramentas e bibliotecas: Python/pandas, R, SPSS, WEKA, etc.}

\section{Aspectos éticos}

\orientacao{Mesmo com dados públicos ou sem participantes humanos diretos, mencione aspectos éticos: anonimização, termos de uso de dados, aprovação em comitê de ética (se necessário).}

\placeholder{Descreva como serão tratados os aspectos éticos da pesquisa. Exemplo: ``Os dados utilizados são públicos e anonimizados, não permitindo identificação de indivíduos.''}
